% STATUS: DRAFT
\chapter{Adding an observer}\label{ch:add_observer}

\begin{physmotiv}
In \cref{ch:static_patch} the static patch algebra is type III$_1$ and $H$ is outside it. In \cref{ch:ds_special} we saw $H$ is a constraint. These two facts are reconciled by including an \emph{observer}: a physical system on the worldline $r=0$ with its own Hamiltonian $q$, whose energy provides a clock. The constraint becomes $H+q=0$; solving it ties the clock reading to the field's modular time and produces, as the algebra of invariants, exactly the crossed product of \cref{ch:crossed}. The sign of $q$, bounded below, is what will later make the trace finite.
\end{physmotiv}

\section{The observer}

Model the observer as a point-like system on the worldline $r=0$ with Hilbert space $\HH_{\rm obs}$ and Hamiltonian $q$ (its energy measured with respect to static time $t$). Physical requirements:
\begin{itemize}
\item \textbf{Energy bounded below:} $q\ge0$ (the observer has positive energy, in particular a rest mass; we absorb its rest mass into the zero of $H$). This is not decoration; it determines the type of the final algebra (\cref{ch:clpw}).
\item \textbf{A clock:} $q$ has a conjugate variable (the observer's proper time, $p$) so that dressed operators ``at time $p$'' can be defined. For an ideal clock, $\HH_{\rm obs}=L^2(\R_{>0})$ or, in the simplest model, $L^2(\R)$ with the projection to $q\ge0$ imposed afterwards.
\end{itemize}
The kinematical Hilbert space is $\HH_{\rm kin}=\HH\otimes\HH_{\rm obs}$ with $\HH$ the QFT Hilbert space in the Bunch--Davies representation.

\section{The constraint}

The static-patch time translation, acting on fields by $H$ and on the observer by $q$, is a gauge symmetry. The gauge-invariant observables commute with the total
\begin{equation}
\hat H=H+q
\end{equation}
(with $H=H_R-H_L$ the two-sided generator for operators in both patches; for operators in $R$ only, $H$ acts as the boost on $\A_R$). Physical states obey $\hat H\Psi=0$ (or, after group averaging, are the group-averaged vectors). The remaining isometries (those that do not preserve the observer's worldline) are broken by its presence; the unbroken ones, rotations about the worldline, form a compact group over which one projects, and we ignore this refinement.

\begin{keyidea}
Physical observables of the static patch with an observer are the operators on $\HH\otimes\HH_{\rm obs}$ commuting with $\hat H=H+q$. Without the observer there are no non-trivial ones localized in the patch.
\end{keyidea}

\section{Dictionary with Part V}

Set $K=\beta_{\rm dS}H$ (the modular Hamiltonian of $\A_R$ in $\Psi_{\rm dS}$) and $x=-\beta_{\rm dS}q$. Then
\begin{equation}
\beta_{\rm dS}\hat H=K-x,
\end{equation}
the constraint generator of \cref{ch:toy_clock} (frame 2). The invariants of $\A_R\otimes\Bnd{L^2(\R_q)}$ under $\Ad\ee^{-\ii t\hat H}$ are generated by the dressed operators
\begin{equation}
\hat a=\ee^{-\ii pK}a\,\ee^{\ii pK}=\sigma_p(a),\quad a\in\A_R;\qquad x=-\beta_{\rm dS}\,q,
\end{equation}
that is, the crossed product $\A_R\rtimes_\sigma\R$ in frame 2. The operator $\hat a$ is the QFT operator $a$ ``evaluated at the observer's clock time'' ($p$ is conjugate to $q$, so shifting by $p$ means translating in the observer's proper time). Physically, $\hat a$ is gauge invariant because the combination (field time $-$ clock time) is what has meaning, not either alone.

\begin{table}[h]
\centering
\begin{tabular}{@{}ll@{}}
\toprule
Crossed product (Part V) & de Sitter with observer\\
\midrule
$\M$ & $\A_R$, type III$_1$ static-patch algebra\\
$\Omega$ & $\Psi_{\rm dS}$ (Bunch--Davies)\\
$K=-\log\Delta$ & $\beta_{\rm dS}H$\\
$x$ & $-\beta_{\rm dS}\,q$ ($q$ = observer energy)\\
constraint $K-x$ & $\beta_{\rm dS}(H+q)$\\
$X=K+x$ (frame 1) & $\beta_{\rm dS}(H-q)$ in the unitarily equivalent frame\\
$x\le0$ & $q\ge0$\\
\bottomrule
\end{tabular}
\end{table}

\section{Why an observer is needed, and what it means algebraically}

Without an observer the static patch algebra is type III$_1$ (\cref{ch:static_patch}) but its elements are not gauge-invariant; with it we get a gauge-invariant algebra of type II. Algebraically, the observer enlarges $\A_R$ by a copy of $\Bnd{L^2}$ (a ``reference frame'') and the constraint cuts it down to the invariants. The result is not a tensor product $\A_R\otimes\A_{\rm obs}$ but a twisted version, because the constraint couples the two. Operationally: the observer \emph{is} the clock relative to which the QFT state is read.

\begin{whatwrong}
\begin{itemize}
\item The assumption that the observer is a point worldline with a decoupled Hamiltonian $q$ is a model. Its mass back-reacts on the geometry (\cref{ch:anatomy}), and its internal structure and finite size matter (\cref{ch:what_observer}).
\item We have not shown that the dressed algebra is \emph{all} of the observables available to the observer; it is the algebra of invariants in the semiclassical ($G_N\to0$) limit.
\item The choice $q\ge0$ and the model of a clock with an unbounded conjugate variable are in tension (a Hamiltonian bounded below has no self-adjoint conjugate time operator, Pauli); we address this in \cref{ch:what_observer}.
\end{itemize}
\end{whatwrong}

\section*{Exercises}
\begin{exercise}
Show that $\sigma_p(a)=\ee^{-\ii pK}a\,\ee^{\ii pK}$ commutes with $K-x$ (use $\ee^{-\ii t(K-x)}\sigma_p(a)\ee^{\ii t(K-x)}$ and $\ee^{\ii tx}p\,\ee^{-\ii tx}=p-t$).
\end{exercise}
\begin{exercise}
For a free scalar and the observer's proper time $p$, write $\hat\phi(f)=\int\dd s\,f(s)\,\ee^{-\ii pK}\phi(s)\ee^{\ii pK}$ for a smeared field along the worldline, and interpret $\hat\phi$ as ``the field measured at clock time $p$''.
\end{exercise}
\begin{exercise}
Explain why the equation $(H+q)\Psi=0$ has no normalizable solutions in $\HH\otimes L^2(\R)$ and how group averaging produces a physical inner product. Compare with \cref{ch:groups}.
\end{exercise}
